% !TeX root = ../常微分方程笔记.tex
% 由主文档通过 \input 载入；请编译主文档。
\setcounter{chapter}{0} % 保留原书章号。
\chapter{数学分析预备知识}
\label{app:ode-analysis}

\section{连续性、紧致性与函数列}
连续函数在紧集上有界且一致连续。对连续向量函数 $u:[a,b]\to\mathbb R^d$，一致范数为 $\|u\|_\infty=\max\|u(t)\|$。若 $u_n$ 在此范数下为 Cauchy 列，则逐点极限 $u$ 存在；对 Cauchy 不等式先令一个指标趋于无穷，可知 $u_n\to u$ 一致。由三角不等式将 $u$ 的局部变化分为两个逼近误差与一个连续函数的变化，便知 $u$ 连续。这证明了 $C([a,b],\mathbb R^d)$ 的完备性。

函数族等度连续（equicontinuous）是指：对每个 $\varepsilon>0$，存在同一个 $\delta>0$，使族中任意 $u$ 和任意 $|s-t|<\delta$ 都满足 $\|u(s)-u(t)\|<\varepsilon$。Arzelà--Ascoli 定理断言，紧区间上一致有界且等度连续的连续函数列存在一致收敛子列。证明的要点是在可数稠密点集上对角抽子列，再用等度连续性把逐点控制扩展为一致控制。有限维值域的有界序列可抽收敛子列，是该论证的重要前提。

\section{多元微分与隐函数定理}
对 $f:\mathbb R^d\to\mathbb R^m$，$Df(x)$ 是由偏导数组成的 $m\times d$ Jacobian 矩阵。当连接 $x,z$ 的线段位于定义域中且 $f\in C^1$ 时，
\[
 f(x)-f(z)=\int_0^1Df(z+\theta(x-z))(x-z)\,\mathrm d\theta.
\]
因此导数有界给出 Lipschitz 估计。多维向量函数一般不能写成对所有分量都共用一个中间点的标量中值公式，积分形式没有这一问题。

隐函数定理的标量版本为：若 $F(t,y)$ 为 $C^1$，$F(t_0,y_0)=0$ 且 $F_y(t_0,y_0)\ne0$，则附近唯一存在 $C^1$ 函数 $y=\phi(t)$ 满足该关系，并有 $\phi'=-F_t/F_y$。对于 $F(t,y,p)=0$ 解出 $p$，使用的非退化条件相应是 $F_p\ne0$。

\section{幂级数与积分中的极限}
幂级数 $\sum a_n(t-t_0)^n$ 在其收敛半径内可逐项求导；在任何严格较小的闭区间上一致收敛。乘积级数与系数比较应在共同的收敛区域内进行，不能只操作形式符号后便宣称得到函数解。

若连续函数 $g_n\to g$ 在有限区间上一致收敛，则
\[
 \left\|\int_a^b g_n-\int_a^b g\right\|\leq(b-a)\|g_n-g\|_\infty\to0.
\]
对无穷区间，若局部一致收敛且 $\|g_n(t)\|\leq h(t)$，其中 $h\geq0$ 可积，则先以 $h$ 控制远端积分，再在有限区间取极限，即可交换极限与积分。交换双重积分时，绝对可积性是本书使用 Fubini 定理的充分条件。对参数积分求导也必须控制导数的可积性或在紧区间上具有连续偏导数，不能只把导数符号移入积分。

